\setcounterpageref{ex}{0}
\TaoPhieuKiemTra{14}

\setcounter{ex}{0}
\setcounter{bt}{0}

\textbf{\large I. PHẦN TRẮC NGHIỆM (3,0 điểm)}

{\itshape Khoanh tròn vào chữ cái đặt trước câu trả lời đúng. Mỗi câu đúng được 0,5 điểm.}

\begin{ex} % Câu 1
Số $68$ được viết thành tổng các chục và đơn vị là
\choice
{$6 + 8$}
{$68 + 0$}
{\True $60 + 8$}
{$60 + 80$}
\loigiai{
$68 = 60 + 8$.\\
\textbf{Đáp án đúng: C}
}
\end{ex}

\begin{ex} % Câu 2
Hiệu của hai số $65$ và $23$ là
\choice
{$88$}
{$32$}
{$43$}
{\True $42$}
\loigiai{
Hiệu: $65 - 23 = 42$.\\
\textbf{Đáp án đúng: D}
}
\end{ex}

\begin{ex} % Câu 3
Một thùng nước có $20\text{ }l$, bác Nam múc ra $8\text{ }l$. Trong thùng còn lại số lít nước là
\choice
{$11\text{ }l$}
{\True $12\text{ }l$}
{$13\text{ }l$}
{$28\text{ }l$}
\loigiai{
Còn lại: $20 - 8 = 12\text{ }l$.\\
\textbf{Đáp án đúng: B}
}
\end{ex}

\begin{ex} % Câu 4
Trên đoạn thẳng $AB$, lấy thêm điểm $M$ nằm giữa $A$ và $B$. Khi đó có tất cả bao nhiêu đoạn thẳng?
\choice
{\True $3$ đoạn thẳng}
{$1$ đoạn thẳng}
{$2$ đoạn thẳng}
{$4$ đoạn thẳng}
\loigiai{
Có 3 đoạn thẳng: $AM, MB, AB$.\\
\textbf{Đáp án đúng: A}
}
\end{ex}

\begin{ex} % Câu 5
Điền dấu thích hợp: $13 - 7 \dots 14 - 8$.
\choice
{$>$}
{$<$}
{Không so sánh được}
{\True $=$}
\loigiai{
$13 - 7 = 6$; $14 - 8 = 6$. Vì $6 = 6$ nên điền dấu $=$.\\
\textbf{Đáp án đúng: D}
}
\end{ex}

\begin{ex} % Câu 6
Tính: $14 - 9 + 7 = \dots$
\choice
{$11$}
{$13$}
{\True $12$}
{$10$}
\loigiai{
$14 - 9 = 5$; $5 + 7 = 12$.\\
\textbf{Đáp án đúng: C}
}
\end{ex}

\vspace{4mm}
\textbf{\large II. PHẦN TỰ LUẬN (7,0 điểm)}

{\itshape Thí sinh trình bày chi tiết bài làm vào giấy kiểm tra.}

\begin{bt}[1,0 điểm]
Tính nhẩm:
\begin{multicols}{2}
\begin{enumerate}[a)]
  \item $8 + 8$
  \item $7 + 9$
  \item $16 - 7$
  \item $18 - 8$
\end{enumerate}
\end{multicols}
\loigiai{
\begin{itemize}
  \item $8 + 8 = 16$
  \item $7 + 9 = 16$
  \item $16 - 7 = 9$
  \item $18 - 8 = 10$
\end{itemize}
}
\end{bt}

\begin{bt}[1,5 điểm]
Đặt tính rồi tính:
\begin{multicols}{3}
\begin{enumerate}[a)]
  \item $51 + 49$
  \item $83 - 23$
  \item $44 + 48$
\end{enumerate}
\end{multicols}
\loigiai{
\begin{itemize}
  \item $51 + 49 = 100$
  \item $83 - 23 = 60$
  \item $44 + 48 = 92$
\end{itemize}
}
\end{bt}

\begin{bt}[1,0 điểm]
Tính giá trị biểu thức kèm đơn vị đo:
\begin{enumerate}[a)]
  \item $30\text{ kg} + 9\text{ kg} - 11\text{ kg}$
  \item $39\text{ }l - 15\text{ }l + 8\text{ }l$
\end{enumerate}
\loigiai{
\begin{enumerate}[a)]
  \item $30\text{ kg} + 9\text{ kg} - 11\text{ kg} = 39\text{ kg} - 11\text{ kg} = 28\text{ kg}$.
  \item $39\text{ }l - 15\text{ }l + 8\text{ }l = 24\text{ }l + 8\text{ }l = 32\text{ }l$.
\end{enumerate}
}
\end{bt}

\begin{bt}[1,5 điểm]
Bạn Châu có $29$ bông hoa. Bạn Dương có ít hơn bạn Châu $6$ bông hoa. Hỏi bạn Dương có bao nhiêu bông hoa?
\loigiai{
Số bông hoa bạn Dương có là:
\[ 29 - 6 = 23 \text{ (bông hoa)} \]
\textbf{Đáp số:} $23$ bông hoa.
}
\end{bt}

\begin{bt}[1,0 điểm]
Một thùng chứa đầy $49\text{ }\text{kg}$ gạo nếp. Buổi sáng cửa hàng đã bán ra $18\text{ }\text{kg}$. Hỏi trong thùng còn lại bao nhiêu ki-lô-gam gạo nếp?
\loigiai{
Trong thùng còn lại số ki-lô-gam gạo nếp là:
\[ 49 - 18 = 31\text{ }\text{kg} \]
\textbf{Đáp số:} $31\text{ }\text{kg}$ gạo nếp.
}
\end{bt}

\begin{bt}[1,0 điểm]
Bạn Nam có một số viên bi. Nếu Nam cho Bình $5$ viên bi thì Nam còn lại $8$ viên bi. Hỏi lúc đầu Nam có bao nhiêu viên bi?
\loigiai{
Lúc đầu Nam có số viên bi là:
\[ 8 + 5 = 13 \text{ (viên bi)} \]
\textbf{Đáp số:} $13$ viên bi.
}
\end{bt}
